\documentclass[10pt]{article}
\usepackage[letterpaper,margin=0.72in]{geometry}
\usepackage{amsmath}
\usepackage{booktabs}
\usepackage{graphicx}
\usepackage{hyperref}
\usepackage{float}
\hypersetup{colorlinks=true,linkcolor=blue,urlcolor=blue}
\setlength{\parindent}{0pt}
\setlength{\parskip}{5pt}
\title{\vspace{-1.2cm}Appendix IV\\Layer Analysis, Ablation, and Propagation Sensitivity}
\author{Sohail Abbas, Muhammad Kazim, Muhammad Fayaz, and Harun Pirim}
\date{}
\begin{document}
\maketitle

\section{Purpose of This Appendix}
This appendix explains why the paper uses four relation layers and how each layer contributes to the final model. It expands the layer-density analysis, aggregation-support figure, ablation table, and propagation-depth sensitivity.

\section{Layer Meaning}
The four relation layers are designed to separate different mechanisms of RPL behavior:
\begin{itemize}
\item \textbf{Routing layer}: parent-child DODAG structure. This is the layer most directly affected by DRA because the attack changes parent-selection incentives.
\item \textbf{Link-quality layer}: weighted version of the observed parent-child support. It distinguishes low-cost and high-cost parent relationships.
\item \textbf{Temporal layer}: similarity in delay, packet-delivery behavior, and route dynamics across observation windows.
\item \textbf{Trust/anomaly layer}: similarity in rank inconsistency, parent-metric deviation, and delivery loss.
\end{itemize}

\section{Layer Structure}
\begin{table}[H]
\centering
\caption{Mean structural properties of the four RPL relation layers.}
\small
\begin{tabular}{lccc}
\toprule
Layer & Mean edges & Density & Edge share\\
\midrule
Routing & 49.00 & 0.040 & 0.023\\
Link quality & 49.00 & 0.040 & 0.023\\
Temporal & 885.78 & 0.723 & 0.410\\
Trust/anomaly & 1154.90 & 0.943 & 0.544\\
\bottomrule
\end{tabular}
\end{table}

The routing and link-quality layers are sparse because each node has a limited parent-child structure. The temporal and trust/anomaly layers are denser because they connect nodes by similarity. If all edges are collapsed into one graph, the dense layers dominate the message-passing support and the model loses the distinction between physical routing relations and behavioral similarity relations.

\section{Aggregation-Support Figure}
\begin{figure}[H]
\centering
\includegraphics[width=0.84\textwidth]{results/cooja_clean_5seed_layer_analysis/figures/fig07_aggregation_support.pdf}
\caption{Aggregation-support analysis. Collapsing relation layers increases graph support and hides relation identity.}
\end{figure}

The figure explains why Agg-GCN is not equivalent to ML-GCN. Agg-GCN receives one combined support, so a message can come from a routing neighbor, a link-quality neighbor, or an anomaly-similar node without preserving which relation produced that edge. ML-GCN keeps the four channels separate and only fuses them after relation-specific convolution.

\section{Layer Ablation}
\begin{table}[H]
\centering
\caption{Single-layer and leave-one-layer-out ML-GCN ablation.}
\small
\begin{tabular}{lccc}
\toprule
Variant & BAcc & F1 & FPR\\
\midrule
All layers (ML-GCN) & \textbf{0.721} & \textbf{0.531} & 0.183\\
Routing only & 0.719 & 0.509 & 0.265\\
Link-quality only & 0.706 & 0.495 & 0.262\\
Without temporal & 0.693 & 0.481 & 0.258\\
Without routing & 0.677 & 0.469 & 0.212\\
Without link quality & 0.671 & 0.471 & \textbf{0.150}\\
Temporal only & 0.525 & 0.337 & 0.801\\
Trust/anomaly only & 0.517 & 0.341 & 0.967\\
\bottomrule
\end{tabular}
\end{table}

The ablation indicates that routing and link-quality information are highly informative for DRA detection, which is expected because decreased-rank attacks manipulate routing attractiveness. However, the full multilayer model obtains the best F1-score, showing that temporal and anomaly information add useful context when combined with the routing layers. The temporal-only and anomaly-only variants perform poorly because behavioral similarity alone is too broad and can connect benign nodes with similar noisy dynamics.

\section{Propagation Sensitivity}
\begin{figure}[H]
\centering
\includegraphics[width=0.92\textwidth]{results/cooja_clean_5seed_graph_sensitivity/figures/fig08_propagation_sensitivity.pdf}
\caption{Propagation-depth sensitivity for graph-based models.}
\end{figure}

The sensitivity experiment compares one-hop and two-hop propagation. In this benchmark, two-hop propagation reduces the detection quality of the graph models. The reason is practical: the LLN snapshots are small, and DRA evidence can be localized around parent-child and nearby behavioral relations. Additional propagation mixes information from wider neighborhoods and can blur the difference between malicious and benign nodes. This supports the paper's choice of one-hop relation-specific message passing.

\end{document}
